\chapter{Keeping the wearer safe}
\label{ch:safety}

The wearer cannot see what the driver intends, cannot feel the controls, and
has a moving machine within reach of their head. This chapter describes what
sits between the two, and, more usefully, the ways that layer can appear to
work while doing nothing.

\section{Three stages and a floor}

\begin{figure}[htbp]
  \centering
  % The layers between a command and a person. Read top to bottom.
\begin{tikzpicture}[x=1mm, y=1mm, font=\scriptsize]

  \node[stage, minimum height=8mm, text width=52mm] (cmd) at (0,0)
    {A wanted hand pose arrives};

  \node[stage, minimum height=9mm, text width=52mm, below=6mm of cmd] (l1)
    {\textbf{1. Solve with obstacles switched on}\\
     rejects poses that would strike something};
  \node[stage, minimum height=9mm, text width=52mm, below=5mm of l1] (l2)
    {\textbf{2. Try again with the elbow moved}\\
     up to six alternative arm shapes for the same hand pose};
  \node[guard, minimum height=9mm, text width=52mm, below=5mm of l2] (l3)
    {\textbf{3. Measure the distance to the wearer}\\
     every moving part, against every body part, from the live positions};
  \node[motor, minimum height=8mm, text width=52mm, below=6mm of l3] (send)
    {Commands go out};

  \draw[go] (cmd) -- (l1); \draw[go] (l1) -- (l2);
  \draw[go] (l2) -- (l3); \draw[go] (l3) -- (send);

  \node[guard, minimum height=22mm, text width=40mm, right=14mm of l2] (stopall)
    {\textbf{Anything here stops the arm}\\[3pt]
     emergency stop button \textbullet\ both master buttons held \textbullet\
     master stops reporting \textbullet\ the safety observer stops confirming
     \textbullet\ any of nine named blocking conditions};
  \draw[halt] (stopall.west) -- (l2.east);
  \draw[halt] (stopall.south) |- (send.east);

  \node[stage, minimum height=15mm, text width=40mm, left=14mm of l2] (see)
    {\textbf{Everything that blocks says so, by name}\\[2pt]
     loud on the first block, louder after three seconds, and a blocker that
     cannot say how it clears is refused at registration};
  \draw[go] (l2.west) -- (see.east);

  \node[tag, text width=100mm, align=left, anchor=north] at ($(send.south)+(0,-8mm)$)
    {Stage 3 is measured from the shapes of the arm and the body directly. It
     does not use the simulator's own collision settings, because those
     deliberately ignore the parts of the arm closest to the wearer's shoulders
     and would report a clear pose with the tube inside the person.};

\end{tikzpicture}

  \caption[The safety layers]{What every command passes through. The first two
  stages reject or reshape a pose. The third measures the distance from the
  arm to the person and holds if it is too small. Anything on the right can
  stop the arm outright, and everything that stops it has to say so by name.}
  \label{fig:safety}
\end{figure}

The first stage asks the solver for joint angles with obstacle checking
switched on, so a pose that would strike the table or the frame comes back
refused. The second stage handles the common case where the hand pose is fine
but the arm has chosen an awkward shape to reach it: the solver is asked again
from up to six different starting guesses, which swings the elbow around
without moving the hand. The third stage is the one that matters, and it is
described next.

\subsection{Why the simulator's own collision check is not enough}

Robot software of this kind ships with a file listing which pairs of parts are
allowed to touch. Pairs that are permanently adjacent, such as the arm's base
and the bracket it is bolted to, are excluded, otherwise every pose would be
reported as a collision.

On a bench robot that is harmless. On this one it is not, because the parts
nearest the wearer are exactly the parts that had to be excluded. Forty-four
such pairs are excluded here, and every one of them is a piece of arm that
swings past a shoulder, a neck or a chest. A pose the standard check calls
valid can have the upper arm tube inside the person.

The clearance check therefore does not use that file. It takes the shapes of
the moving arm and the shapes of the body, from their live positions, and
computes the shortest distance between them. If that distance falls below
\SI{150}{\milli\metre}, the command is held, nothing is published, and the log
names which body part was too close.

Two facts about that number are worth stating. The mounting bracket itself
sits \SI{161}{\milli\metre} from the wearer's torso and no joint moves it, so
the margin has only \SI{11}{\milli\metre} of headroom before the arm has done
anything at all. And a clearance reading of exactly \SI{161}{\milli\metre} is
a measurement of the bracket rather than of the arm, which is why so many
figures in this project sit at that value.

\section{Where the body is}

Until late in the project the body in the model was a mannequin, and the
clearance check measured against that. A camera was then added that watches
the wearer and estimates where their limbs actually are, using an off-the-shelf
body-tracking model~\cite{lugaresi2019mediapipe}. Merging a measured body with
an assumed one is where this could go wrong, so the rule is deliberately
one-sided.

\begin{quote}
The camera may only ever make the wearer \emph{bigger}.
\end{quote}

Part by part, whichever of the measured shape and the assumed shape is
\emph{closer} to the robot is the one the check uses. A tracking result placing
the body further away than assumed changes nothing at all. This was tested by
injecting a body \SI{0.1}{\metre}, \SI{0.3}{\metre}, \SI{1}{\metre} and
\SI{5}{\metre} further away in turn; none of them moved the enforced margin.
Six deliberate faults were then fed through the running system, including a
dark frame, an empty frame and two people in shot, and each refusal was
required to name the fault it had caught. That last requirement is what
exposed a check that had never actually run.

The honest limitation is on the other side. If the wearer raises their arms
and the tracking then drops out, the system falls back to the assumed posture,
which has the arms down. Measured through the guard's own geometry, arms held
out to the sides leave \SI{93}{\milli\metre} of clearance at the home pose,
where the fallback believes there is \SI{317}{\milli\metre}. The current answer
to that is an instruction to the wearer, not a control system, and
\cref{ch:discussion} treats it as the open safety problem it is.

\section{Stopping}

Three things stop the arms, and they are independent by design: a button, both
master buttons held together, and the master going quiet.

The third is the one that taught the most. A dead-man switch of this kind
watches for the driver's input to stop arriving. The implementation watched
for messages to stop arriving, which sounds like the same thing and is not.
When the master's sensors degrade, the software that reads them does not go
silent; it substitutes the last good reading and carries on publishing at the
same rate. The topic never went quiet, so the dead-man could never fire on the
failure that actually happens. It fired reliably when the cable was pulled out,
which is why it passed every test for months.

The fix was to stamp each frame with the age of the data it was built from
rather than the time it was sent, and to key the dead-man on that. Measured
afterwards: a master publishing steadily but with frozen data trips in
\SIrange{0.94}{0.99}{\second}; a genuinely silent master trips in
\SI{0.97}{\second}; and 243 healthy frames produce no false trip.

A second fault in the same mechanism is worth recording because it was
invisible in the same way. The stop routine tried to notify the arm driver
first, using a call that waits up to two seconds for a service that does not
exist when running in simulation. It did that twice. So every emergency stop
blocked for four seconds before parking the arms, before logging anything, and
before the rest of the system could be told. The detection was always on time.
The response was not, and the log line that would have revealed it was itself
stuck behind the delay. Parking now happens first and the driver is notified
afterwards.

\section{Making blocking visible}

Four separate mechanisms in this system can stop all motion, and each of them
is usually right to. The problem was never the stopping. It was that a stopped
arm and a working arm looked identical from outside.

Everything that can stop motion now registers with one component that enforces
three rules. It must announce itself by name the first time it blocks and
every five seconds after. Blocking for more than three seconds is escalated,
because sustained blocking is never normal. And a blocker has to state, at the
moment it registers, how it will clear; one that cannot is rejected outright,
because a condition with no exit is a latch.

That last rule came out of an incident which repeated five times, the fifth
inside the code written to catch the first four. The pattern is always the
same: the code that clears a condition sits after the code that sets it, so if
setting the condition also stops that routine running, the condition can never
be cleared. Rather than police each instance, any blocker that stops being
re-asserted for two seconds now expires, and an expiry is treated as a third
state rather than as an all-clear. Nobody asserting a condition is not evidence
that the condition went away; it is evidence that whatever was watching has
stopped.

\section{Testing by breaking things}

The safety layer is validated by injecting faults rather than by reading the
code. Fourteen faults are injected in turn, covering sensor dropout, the
master disconnecting, an emergency stop during motion, the object-tracking
losing sight of its target, a logging failure, a clock going backwards, an
attempt to write identifying data, the arm's control session being lost and
the network to an arm going down. All fourteen are handled.

Getting there was less tidy than that sentence suggests. Four separate bugs in
the injection harness each produced a clean table of results while testing
nothing at all. A fixed delay was not reliably landing inside a trial. The
publisher had not finished being discovered before the fault was sent. State
left over from the previous run satisfied the trigger condition for the next
one. And the abort signal was a one-shot event, so a fault raised in the gap
between two trials was seen by neither.

The transferable lesson is not about this system. Any fault-injection harness
is itself an instrument, and its verdicts should not be believed until it has
been shown capable of reporting a failure.

\begin{table}[htbp]
  \centering
  \small
  \caption[Fault injection]{The fourteen injected faults and the response
  required of each. All are currently handled.}
  \label{tab:faults}
  \begin{tabular}{@{}p{0.34\linewidth}p{0.58\linewidth}@{}}
    \toprule
    \textbf{Injected fault} & \textbf{Required response} \\
    \midrule
    A sensor channel drops out & Channel disabled automatically, loss of
      capability announced, trial marked invalid \\
    The master is unplugged & Dead-man trips within one second, arms park \\
    Emergency stop during motion & Arms park, condition latches, clears only
      on an explicit reset \\
    Object tracking latches the wrong target & Estimate switches within
      \SI{0.1}{\second} \\
    Two targets equally likely & No assistance offered; a tie is not a
      decision \\
    The tracked object disappears & Estimate cleared rather than held \\
    Nothing is detected at all & Assistance yields to the driver \\
    The log cannot be written & Raises rather than dropping samples \\
    Identifying data is submitted & Refused in code \\
    The clock goes backwards & No interval measured from wall-clock time \\
    The master reconnects & Detected again, next trial gated until data is
      fresh \\
    The arm's control session is lost & Closed and reopened without a restart \\
    The camera sees nothing before a trial & Trial refused before it starts \\
    The network to one arm drops & \emph{Both} arms frozen \\
    \bottomrule
  \end{tabular}
\end{table}

The last row is a design decision rather than a technical one. One arm live
and one arm frozen on a person's back is worse than both frozen, and the two
cannot be told apart by looking.

\section{Two problems the desk arrangement introduces}
\label{sec:tracking}

When the driver holds tracked hand controllers across the room, the headset
sits on a shelf and acts as the fixed reference for everything the controllers
report. Two failures follow from that and neither announces itself.

The first is that the shelf can be knocked. Nothing breaks: the controllers
stay tracked, the update rate stays at \SI{90}{\hertz}, every reported position
remains valid, and all of them are now expressed in a frame that has quietly
rotated. The system latches the reference at the start and freezes if it moves
by more than \SI{20}{\milli\metre} or \SI{2}{\degree}. The rotation limit is the
one that matters, because a small twist barely moves the headset and puts every
subsequent command several degrees wrong.

The second is that a controller can be blocked from view. The system had a
signal named for exactly this and it was derived from when the last message
arrived, not from whether the controller in it was being tracked. Covering a
controller with a hand for five seconds produced no reaction at all while the
stream ran on normally. When the driver wears the headset this hardly matters,
because their hands are in front of its cameras. With the headset on a shelf,
being blocked from view is routine.

\begin{figure}[htbp]
  \centering
  \screenshot{58mm}{The safety view during a run: the wearer drawn with the
  \SI{150}{\milli\metre} margin shown as a shell around them, the nearest part
  of the arm highlighted and labelled with which body part it is closest to,
  and the readout showing the current clearance. Take one shot in a safe pose
  and one with the arm deliberately brought close, so the difference is
  visible.}
  \caption[The safety view]{The clearance view. It draws the same body the
  check is enforcing against, so what is displayed and what is enforced cannot
  drift apart.}
  \label{fig:safetyview}
\end{figure}
